% ============================================================================
\chapter{椭圆曲线初步}
\label{ch:elliptic-curves-intro}
% ============================================================================

本章介绍椭圆曲线的基本概念，为后续将模形式与椭圆曲线联系起来做准备。
椭圆曲线是代数几何中最简单的非平凡对象——亏格为 $1$ 的代数曲线——但其算术
性质极其深刻。费马大定理的证明正是建立在椭圆曲线与模形式之间的桥梁之上。

% ============================================================================
\section{Weierstrass 方程}
\label{sec:weierstrass-equation}
% ============================================================================

\begin{definition}[椭圆曲线]
\label{def:elliptic-curve}
设 $K$ 为域。$K$ 上的一条\textbf{椭圆曲线}（elliptic curve）是一条亏格为
$1$ 的光滑射影曲线 $E$，并指定一个 $K$-有理点 $O \in E(K)$ 作为
\textbf{单位元}。
\end{definition}

\begin{definition}[Weierstrass 方程]
\label{def:weierstrass-equation}
椭圆曲线的\textbf{一般 Weierstrass 方程}为
\[
  y^2 + a_1 xy + a_3 y = x^3 + a_2 x^2 + a_4 x + a_6,
  \quad a_i \in K.
\]
当 $\mathrm{char}(K) \neq 2, 3$ 时，通过变量替换可化简为
\textbf{短 Weierstrass 方程}：
\[
  y^2 = x^3 + ax + b, \quad a, b \in K.
\]
\end{definition}

\begin{definition}[判别式与 $j$-不变量]
\label{def:discriminant-j}
短 Weierstrass 方程 $y^2 = x^3 + ax + b$ 的\textbf{判别式}为
\[
  \Delta = -16(4a^3 + 27b^2).
\]
曲线光滑（即为椭圆曲线）当且仅当 $\Delta \neq 0$。

\textbf{$j$-不变量}定义为
\[
  j(E) = -1728 \frac{(4a)^3}{\Delta} = \frac{-1728 \cdot 64 a^3}{-16(4a^3 + 27b^2)}
  = \frac{6912 a^3}{16(4a^3 + 27b^2)}.
\]
更简洁地写为 $j = 1728 \cdot \frac{4a^3}{4a^3 + 27b^2}$。

两条椭圆曲线在 $\bar{K}$ 上同构当且仅当它们有相同的 $j$-不变量。
\end{definition}

\begin{proposition}[$\Delta \neq 0$ 的几何意义]
\label{prop:discriminant-smooth}
$\Delta = 0$ 当且仅当曲线 $y^2 = x^3 + ax + b$ 有奇异点（尖点或结点）。
\end{proposition}

\begin{proof}
曲线 $C\colon F(x,y) = y^2 - x^3 - ax - b = 0$ 在点 $(x_0, y_0)$ 处奇异，
当且仅当
\[
  F = 0, \quad \frac{\partial F}{\partial x} = -3x_0^2 - a = 0,
  \quad \frac{\partial F}{\partial y} = 2y_0 = 0.
\]
由 $y_0 = 0$ 和 $F = 0$ 得 $x_0^3 + ax_0 + b = 0$，即 $x_0$ 是
$f(x) = x^3 + ax + b$ 的根。由 $3x_0^2 + a = 0$ 得 $f'(x_0) = 0$，
即 $x_0$ 是 $f$ 的重根。

$f$ 有重根当且仅当 $\gcd(f, f') \neq 1$。计算 $f(x) = x^3 + ax + b$，
$f'(x) = 3x^2 + a$。用辗转相除法：
\[
  f(x) = \frac{x}{3} f'(x) + \frac{2a}{3}x + b,
\]
\[
  f'(x) = \frac{3x}{2a}\left(\frac{2a}{3}x + b\right)
  + a - \frac{3b^2}{4a} \cdot \frac{3}{2a}
\]
（当 $a \neq 0$）。最终，$f$ 有重根等价于判别式
$\mathrm{disc}(f) = -(4a^3 + 27b^2) = 0$。

当 $a = 0$ 时，$f = x^3 + b$，有重根当且仅当 $b = 0$，
此时 $4a^3 + 27b^2 = 0$，一致。
\end{proof}

\begin{example}
\label{ex:elliptic-curves-examples}
\leavevmode
\begin{enumerate}
  \item $E\colon y^2 = x^3 - x$。此时 $a = -1$，$b = 0$，
    $\Delta = -16(4(-1)^3 + 0) = 64 \neq 0$。
    $j = 1728 \cdot \frac{-4}{-4} = 1728$。
  \item $E\colon y^2 = x^3 + 1$。此时 $a = 0$，$b = 1$，
    $\Delta = -16 \cdot 27 = -432 \neq 0$。$j = 0$。
  \item $y^2 = x^3$（尖点奇异曲线）：$\Delta = 0$，不是椭圆曲线。
  \item $y^2 = x^3 - x^2$（结点奇异曲线）：可化为短 Weierstrass 形式后
    验证 $\Delta = 0$。
\end{enumerate}
\end{example}

\begin{remark}[与复环面的联系]
\label{rem:ec-torus}
在 $\C$ 上，每条椭圆曲线 $E$ 都同构于一个复环面 $\C/\Lambda$
（\cref{ex:riemann-surface-examples}(3)）。同构由 Weierstrass $\wp$ 函数
（\cref{def:weierstrass-p}）给出：
\[
  \C/\Lambda \xrightarrow{\;\sim\;} E(\C), \quad
  z \mapsto [\wp(z) : \wp'(z) : 1], \quad 0 \mapsto O = [0:1:0].
\]
这里 $E$ 的方程正是 $(\wp')^2 = 4\wp^3 - g_2\wp - g_3$
（\cref{thm:weierstrass-properties}(4)）。
\end{remark}


% ============================================================================
\section{群结构}
\label{sec:group-law}
% ============================================================================

椭圆曲线上有自然的 Abel 群结构，单位元为无穷远点 $O = [0:1:0]$。

\begin{theorem}[椭圆曲线的群律]
\label{thm:group-law}
设 $E\colon y^2 = x^3 + ax + b$ 为 $K$ 上的椭圆曲线。
$E(K)$ 上有唯一的 Abel 群结构，以 $O$ 为单位元，满足：
三个点 $P, Q, R \in E$ 共线当且仅当 $P + Q + R = O$。
\end{theorem}

\begin{proof}[构造与验证]
我们显式给出加法公式。设 $P = (x_1, y_1)$，$Q = (x_2, y_2) \in E(K)$，
$P, Q \neq O$。

\textbf{情形 1：$P \neq Q$，$x_1 \neq x_2$}（弦法）。

过 $P, Q$ 的直线斜率为 $\lambda = \frac{y_2 - y_1}{x_2 - x_1}$，
直线方程 $y = \lambda(x - x_1) + y_1$。代入 $y^2 = x^3 + ax + b$，
得到关于 $x$ 的三次方程：
\[
  [\lambda(x - x_1) + y_1]^2 = x^3 + ax + b.
\]
展开整理为 $x^3 - \lambda^2 x^2 + \cdots = 0$。
由 Vieta 定理，三个根 $x_1, x_2, x_3$ 满足
$x_1 + x_2 + x_3 = \lambda^2$，故第三个交点的坐标为
\begin{align}
  x_3 &= \lambda^2 - x_1 - x_2, \\
  y_3 &= \lambda(x_1 - x_3) - y_1.
\end{align}
定义 $P + Q = (x_3, y_3)$（即第三交点关于 $x$ 轴的对称点）。

\textbf{情形 2：$P = Q$，$y_1 \neq 0$}（切法）。

对曲线方程隐函数求导：
$2y\frac{dy}{dx} = 3x^2 + a$，故切线斜率
$\lambda = \frac{3x_1^2 + a}{2y_1}$。
然后用与情形 1 相同的公式计算 $2P = P + P$。

\textbf{情形 3：$P = Q$，$y_1 = 0$}。

此时切线垂直（$x = x_1$），与 $E$ 的第三交点是 $O$。故 $2P = O$，
即 $P$ 是 $2$-挠点。

\textbf{情形 4：$x_1 = x_2$，$y_1 = -y_2$}（$P$ 与 $Q$ 关于 $x$ 轴对称）。

连线垂直，与 $E$ 交于 $O$。故 $P + Q = O$，即 $Q = -P$。

\textbf{逆元}：$-P = (x_1, -y_1)$。

\textbf{验证结合律}：直接代数验证非常繁琐但可行。更优雅的证明利用
除子（divisor）理论：在 $E$ 上定义群律为
$P + Q = R$ 当且仅当 $(P) + (Q) - (R) - (O) = \mathrm{div}(f)$
对某有理函数 $f$ 成立。结合律由除子加法的结合律保证。

具体地，定义映射 $\kappa\colon E(K) \to \mathrm{Pic}^0(E)$，
$P \mapsto [(P) - (O)]$。这里 $\mathrm{Pic}^0(E)$ 是度为零的除子类群，
其群结构是标准的。可以证明 $\kappa$ 是群同构（利用 Riemann-Roch 定理
证明满射：对任意度为零的除子 $D$，$D + (O)$ 是度为 $1$ 的除子，
由 $g = 1$ 的 Riemann-Roch，$\ell(D + (O)) = 1$，故存在唯一的
$P \in E$ 使得 $D \sim (P) - (O)$）。
\end{proof}

\begin{example}
\label{ex:group-law-computation}
在 $E\colon y^2 = x^3 - 5x + 8$ 上，设 $P = (1, 2)$，$Q = (3, 4)$。

斜率 $\lambda = \frac{4 - 2}{3 - 1} = 1$。

$x_3 = 1^2 - 1 - 3 = -3$，
$y_3 = 1 \cdot (1 - (-3)) - 2 = 2$。

故 $P + Q = (-3, 2)$。验证：$(-3)^3 - 5(-3) + 8 = -27 + 15 + 8 = -4$，
但 $y^2 = 4$，即 $2^2 = 4$。正确。
\end{example}


% ============================================================================
\section{同源}
\label{sec:isogeny}
% ============================================================================

\begin{definition}[同源]
\label{def:isogeny}
设 $E_1, E_2$ 为 $K$ 上的椭圆曲线。一个\textbf{同源}（isogeny）
$\varphi\colon E_1 \to E_2$ 是一个满足 $\varphi(O_1) = O_2$ 的
代数群态射（即既是代数簇的态射，又是群同态）。
\end{definition}

\begin{proposition}
\label{prop:isogeny-properties}
设 $\varphi\colon E_1 \to E_2$ 为非零同源。则：
\begin{enumerate}
  \item $\varphi$ 是满射（作为 $\bar{K}$-点集之间的映射）。
  \item $\ker \varphi = \varphi^{-1}(O_2)$ 是有限群。
  \item 存在唯一的正整数 $\deg \varphi$（称为\textbf{度}），等于函数域扩张
    $[K(E_1) : \varphi^* K(E_2)]$ 的度。
  \item $\#\ker \varphi \leq \deg \varphi$，等号在\textbf{可分}同源
    （separable isogeny）时成立。
\end{enumerate}
\end{proposition}

\begin{proof}
\textbf{(1)}：非常值全纯映射（或代数簇的有限态射）在紧（射影）曲线间必为满射。
这是代数几何的标准事实：对射影曲线 $C_1 \to C_2$ 的态射，
像是闭子集（射影簇的像在射影簇中闭），且若非常值则像不是单点，
在不可约曲线上只能是全部 $C_2$。

\textbf{(2)}：$\ker \varphi$ 是 $E_1$ 的闭子群。若 $\ker \varphi$ 无限，
则 $\ker \varphi$ 包含 $E_1$ 的某个不可约分支（不可约曲线的无限闭子集只能是全部），
即 $\ker \varphi = E_1$，但这意味着 $\varphi$ 把 $E_1$ 映到单点 $O_2$，
矛盾于 $\varphi$ 非零。

\textbf{(3)}：$\varphi$ 诱导函数域嵌入 $\varphi^*\colon K(E_2) \hookrightarrow K(E_1)$。
度定义为此域扩张的度。这是有限扩张，因为 $E_1, E_2$ 同维
（都是曲线）且 $\varphi$ 满射。

\textbf{(4)}：在可分情形下，$\deg \varphi$ 等于一般纤维的元素个数，
即 $\#\varphi^{-1}(Q)$ 对一般的 $Q$ 相同，等于 $\deg \varphi$。
而 $\ker \varphi = \varphi^{-1}(O_2)$ 是一条特殊纤维，
由群同态性质，所有纤维有相同的大小。
非可分情形下纤维可能更小（因为有"重合"的点）。
\end{proof}

\begin{definition}[对偶同源]
\label{def:dual-isogeny}
设 $\varphi\colon E_1 \to E_2$ 为度 $n$ 的同源。存在唯一的同源
$\hat{\varphi}\colon E_2 \to E_1$（称为 $\varphi$ 的\textbf{对偶同源}，
dual isogeny），满足
\[
  \hat{\varphi} \circ \varphi = [n]_{E_1}, \quad
  \varphi \circ \hat{\varphi} = [n]_{E_2},
\]
其中 $[n]$ 表示乘以 $n$ 的映射（即 $[n](P) = P + P + \cdots + P$，
$n$ 次）。且 $\deg \hat{\varphi} = \deg \varphi = n$。
\end{definition}

\begin{example}[乘法映射]
\label{ex:multiplication-map}
$[n]\colon E \to E$ 是度 $n^2$ 的同源。其对偶是自身：$\widehat{[n]} = [n]$。
\end{example}


% ============================================================================
\section{挠点与 Tate 模}
\label{sec:tate-module}
% ============================================================================

\begin{definition}[$n$-挠子群]
\label{def:torsion}
椭圆曲线 $E$ 的 \textbf{$n$-挠子群}定义为
\[
  E[n] = \ker([n]) = \{P \in E(\bar{K}) \mid nP = O\}.
\]
\end{definition}

\begin{theorem}[$n$-挠子群的结构]
\label{thm:torsion-structure}
设 $E$ 为域 $K$ 上的椭圆曲线。
\begin{enumerate}
  \item 若 $\mathrm{char}(K) = 0$ 或 $\mathrm{char}(K) \nmid n$，则
    $E[n] \cong (\Z/n\Z)^2$。
  \item 若 $\mathrm{char}(K) = p > 0$，则
    $E[p^e] \cong (\Z/p^e\Z)$ 或 $E[p^e] = \{O\}$，
    取决于 $E$ 是\textbf{寻常}（ordinary）还是\textbf{超奇异}（supersingular）的。
\end{enumerate}
\end{theorem}

\begin{proof}
\textbf{(1)}：$[n]\colon E \to E$ 的度为 $n^2$（这可由复分析证明：
在 $\C/\Lambda$ 上，$[n]$ 将 $\Lambda$ 映到 $n\Lambda$，
$\deg[n] = [\Lambda : n\Lambda] = n^2$）。
当 $\gcd(n, \mathrm{char}(K)) = 1$ 时，$[n]$ 是可分的
（这可由形式群的理论证明：$[n]$ 的形式导数是 $n$，
在 $\gcd(n, p) = 1$ 时非零）。
由\cref{prop:isogeny-properties}(4)，$\#E[n] = \deg[n] = n^2$。

$E[n]$ 是交换群，且 $n \cdot P = O$ 对所有 $P \in E[n]$，故
$E[n]$ 是 $(\Z/n\Z)$-模。有限交换群的结构定理给出
$E[n] \cong \bigoplus_{i} \Z/d_i\Z$（$d_i | n$）。
由 $\#E[n] = n^2$，我们需要 $\prod d_i = n^2$。

关键约束来自 Weil 配对（见下文\cref{def:weil-pairing}）：存在
非退化的双线性配对 $e_n\colon E[n] \times E[n] \to \mu_n$。
非退化性要求 $E[n]$ 中存在阶为 $n$ 的元素。
由此可以推出 $E[n] \cong (\Z/n\Z)^2$。

\textbf{(2)}：当 $\mathrm{char}(K) = p$ 时，$[p]$ 不可分
（形式导数为 $p = 0$），故 $\#E[p] < \deg[p] = p^2$。
根据 $[p]$ 的可分度，有 $\#E[p] = p$（寻常情形）或
$\#E[p] = 1$（超奇异情形）。
\end{proof}

\begin{definition}[Weil 配对]
\label{def:weil-pairing}
对正整数 $n$（与 $\mathrm{char}(K)$ 互素），存在双线性映射
\[
  e_n\colon E[n] \times E[n] \to \mu_n,
\]
称为 \textbf{Weil 配对}（Weil pairing），其中 $\mu_n$ 是 $n$ 次单位根群。
它满足：
\begin{enumerate}
  \item 双线性：$e_n(P + P', Q) = e_n(P, Q) e_n(P', Q)$，
    $e_n(P, Q + Q') = e_n(P, Q) e_n(P, Q')$；
  \item 反对称：$e_n(P, Q) = e_n(Q, P)^{-1}$；
  \item 非退化：若 $e_n(P, Q) = 1$ 对所有 $Q \in E[n]$，则 $P = O$；
  \item Galois 等变：$e_n(\sigma P, \sigma Q) = \sigma(e_n(P, Q))$
    对所有 $\sigma \in G_K$。
\end{enumerate}
\end{definition}

\begin{definition}[Tate 模]
\label{def:tate-module}
设 $\ell$ 为素数（$\ell \neq \mathrm{char}(K)$）。$E$ 的
\textbf{$\ell$-进 Tate 模}（Tate module）定义为逆极限
\[
  T_\ell(E) = \varprojlim_{n} E[\ell^n],
\]
其中转移映射为 $[\ell]\colon E[\ell^{n+1}] \to E[\ell^n]$。

由\cref{thm:torsion-structure}，$T_\ell(E) \cong \Z_\ell^2$
作为 $\Z_\ell$-模。定义
\[
  V_\ell(E) = T_\ell(E) \otimes_{\Z_\ell} \Q_\ell \cong \Q_\ell^2.
\]
\end{definition}

\begin{remark}
\label{rem:tate-galois}
绝对 Galois 群 $G_K = \Gal(\bar{K}/K)$ 自然作用在 $E[\ell^n]$ 上
（因为 $\ell^n$-挠点由代数方程定义），由此诱导连续表示
\[
  \rho_{E,\ell}\colon G_K \to \GL(T_\ell(E)) \cong \GL_2(\Z_\ell),
\]
即一个 $2$ 维的 $\ell$-进 \textbf{Galois 表示}。
这个表示将在\cref{ch:galois-representations}中详细讨论，
它是连接椭圆曲线与模形式的关键桥梁。
\end{remark}


% ============================================================================
\section{$\Q$ 上椭圆曲线的约化}
\label{sec:reduction}
% ============================================================================

\begin{definition}[极小 Weierstrass 方程]
\label{def:minimal-weierstrass}
设 $E/\Q$ 为椭圆曲线。对每个素数 $p$，$E$ 有一个在 $\Z_p$ 上的
\textbf{极小 Weierstrass 方程}——即系数在 $\Z_p$ 中，且判别式的
$p$-进赋值 $v_p(\Delta)$ 最小的 Weierstrass 方程。
（存在全局极小方程，即对所有 $p$ 同时极小。）
\end{definition}

\begin{definition}[约化类型]
\label{def:reduction-type}
设 $E/\Q$ 有极小 Weierstrass 方程，将其系数模 $p$ 约化得到
$\F_p$ 上的曲线 $\tilde{E}$。

\begin{enumerate}
  \item \textbf{好约化}（good reduction）：$\tilde{E}$ 仍是椭圆曲线
    （$p \nmid \Delta$）。此时定义
    \[
      a_p(E) = p + 1 - \#\tilde{E}(\F_p).
    \]
  \item \textbf{乘性约化}（multiplicative reduction）：$\tilde{E}$ 有一个结点。
    \begin{itemize}
      \item \textbf{分裂乘性约化}（split）：切线斜率在 $\F_p$ 中，$a_p = 1$；
      \item \textbf{非分裂乘性约化}（non-split）：切线斜率在 $\F_{p^2} \setminus \F_p$ 中，$a_p = -1$。
    \end{itemize}
  \item \textbf{加性约化}（additive reduction）：$\tilde{E}$ 有一个尖点，$a_p = 0$。
\end{enumerate}
\end{definition}

\begin{theorem}[Hasse 界]
\label{thm:hasse-bound}
设 $E/\Q$ 在素数 $p$ 处有好约化。则
\[
  |a_p(E)| \leq 2\sqrt{p}.
\]
\end{theorem}

\begin{proof}
我们给出证明的核心思路。设 $\phi\colon \tilde{E} \to \tilde{E}$ 为
Frobenius 自同态 $\phi(x, y) = (x^p, y^p)$。则
$\#\tilde{E}(\F_p) = \deg(1 - \phi)$（因为 $\F_p$-有理点恰好是
$1 - \phi$ 的不动点）。

在 $T_\ell(\tilde{E})$ 上，$\phi$ 的特征多项式为
$t^2 - a_p t + p$（由 Eichler-Shimura 关系或直接计算），
其根 $\alpha, \beta$ 满足 $\alpha + \beta = a_p$，$\alpha\beta = p$。

关键不等式：由 $\deg(m - n\phi) \geq 0$ 对所有 $m, n \in \Z$
（度是非负的），利用 $\deg(m - n\phi) = m^2 - a_p mn + pn^2$，
这是关于 $m/n$ 的二次函数非负，要求判别式
$a_p^2 - 4p \leq 0$，即 $|a_p| \leq 2\sqrt{p}$。
\end{proof}

\begin{definition}[导子]
\label{def:conductor}
$E/\Q$ 的\textbf{导子}（conductor）$N_E$ 定义为
\[
  N_E = \prod_{p \text{ 坏约化}} p^{f_p},
\]
其中指数 $f_p \geq 1$ 由 $E$ 在 $p$ 处的约化类型决定：
\begin{itemize}
  \item 乘性约化：$f_p = 1$；
  \item 加性约化：$f_p = 2 + \delta_p$，其中 $\delta_p \geq 0$
    是"野蛮部分"（wild part），仅当 $p = 2$ 或 $3$ 时可能非零。
\end{itemize}
导子 $N_E$ 编码了 $E$ 在哪些素数处有坏约化以及坏约化的程度。
\end{definition}

\begin{remark}[半稳定椭圆曲线]
\label{rem:semistable}
若 $E/\Q$ 在每个素数 $p$ 处要么有好约化，要么有乘性约化
（即没有加性约化），则称 $E$ 是\textbf{半稳定}（semistable）的。
此时 $N_E$ 是无平方因子的（squarefree）。
Wiles 在 1995 年首先证明的是半稳定椭圆曲线的模性
（\cref{ch:wiles-proof}），这足以推出费马大定理。
\end{remark}


% ============================================================================
\section{椭圆曲线的 $L$-函数}
\label{sec:ec-l-function}
% ============================================================================

\begin{definition}[Hasse-Weil $L$-函数]
\label{def:hasse-weil-l}
设 $E/\Q$ 为椭圆曲线，导子为 $N$。定义 $E$ 的 \textbf{Hasse-Weil $L$-函数}为
Euler 乘积
\[
  L(E, s) = \prod_{p \nmid N} \frac{1}{1 - a_p p^{-s} + p^{1-2s}}
  \cdot \prod_{p \mid N} \frac{1}{1 - a_p p^{-s}},
\]
其中 $a_p$ 如\cref{def:reduction-type}中所定义。

等价地，写成 Dirichlet 级数
\[
  L(E, s) = \sum_{n=1}^{\infty} \frac{a_n}{n^s},
\]
其中系数 $a_n$ 由 Euler 乘积的展开确定（乘性函数）。
由 Hasse 界（\cref{thm:hasse-bound}），此 Dirichlet 级数在
$\Re(s) > 3/2$ 时绝对收敛。
\end{definition}

\begin{remark}[模性猜想的 $L$-函数版本]
\label{rem:modularity-l-function}
模性猜想（\cref{conj:modularity}）断言：$L(E, s) = L(f, s)$ 对某个
权 $2$ 新形式 $f \in \Sk_2^{\mathrm{new}}(\GammaO(N_E))$ 成立。
这意味着 $L(E, s)$ 可以解析延拓到整个复平面并满足函数方程——
这些性质单从椭圆曲线的定义出发极难证明，但从模形式的角度则是自然的。
\end{remark}


% ============================================================================
\section{习题}
\label{sec:ch2-exercises}
% ============================================================================

\begin{exercise}
\label{ex:ch2-1}
验证 $E\colon y^2 = x^3 - 43x + 166$ 上的点 $P = (3, 8)$ 和
$Q = (-5, 16)$ 确实在曲线上，并计算 $P + Q$ 和 $2P$。
\end{exercise}

\begin{exercise}
\label{ex:ch2-2}
设 $E\colon y^2 = x^3 + 1$ 在 $\F_7$ 上。
列出 $E(\F_7)$ 的所有点，计算 $a_7(E)$，并验证 Hasse 界。
\end{exercise}

\begin{exercise}
\label{ex:ch2-3}
证明：若 $E/\Q$ 有 $j$-不变量 $j = 0$，则 $E$ 同构于
$y^2 = x^3 + b$ 的形式（$b \in \Q^*$）。
类似地，若 $j = 1728$，则 $E$ 同构于 $y^2 = x^3 + ax$（$a \in \Q^*$）。
\end{exercise}

\begin{exercise}
\label{ex:ch2-4}
设 $\varphi\colon E_1 \to E_2$ 为度 $n$ 的同源。
证明 $\hat{\hat{\varphi}} = \varphi$，且 $\deg(\varphi + \psi) \leq
(\sqrt{\deg \varphi} + \sqrt{\deg \psi})^2$ 对同源 $\varphi, \psi\colon E_1 \to E_2$。
\end{exercise}

\begin{exercise}
\label{ex:ch2-5}
设 $E\colon y^2 = x^3 - x$ 在 $\F_p$ 上（$p$ 为奇素数）。
证明：
\begin{enumerate}
  \item 若 $p \equiv 3 \pmod{4}$，则 $a_p = 0$（即 $\#E(\F_p) = p + 1$）。
  \item 若 $p \equiv 1 \pmod{4}$，写 $p = a^2 + b^2$（$a$ 为奇数，$b$ 为偶数，
    $a > 0$），则 $a_p = \pm 2a$。
\end{enumerate}
（提示：对 (1)，利用 $-1$ 不是模 $p$ 的二次剩余。）
\end{exercise}
